\chapter{The scalar Parisi problem on the support}\label{sec:scalar}

In this section, we collect the properties of the scalar Parisi equation used
in the barrier argument. Recall from \cref{sec:inputs} the solution
$\Phi_\alpha$ of the Parisi equation for a distribution function $\alpha$,
its spatial derivatives $u$, $\chi$, and $w$, the scalar diffusion $Y_t$,
and the function $\Gamma(t)=\E u(t,Y_t)^2$. We first prove two statements
that hold for every distribution function. In \cref{lem:sc-derivatives}, we
bound each positive spatial derivative of $\Phi_\alpha$ by a constant that
depends only on $\beta$ and on the order of the derivative. In
\cref{lem:sc-ito}, we record the stochastic differential equations that these
derivatives satisfy along $Y$. We then specialize to the Parisi distribution
function $\alpha_\beta$ on the interval $(0,q_\beta)$, where $\Gamma(t)=t$.
In \cref{b:lem:os-scalar-jets}, we prove the two stationarity identities
$\beta^2\E\chi(t,Y_t)^2=1$ and $\E w(t,Y_t)^2=2\alpha(t)\E\chi(t,Y_t)^3$,
the formula $\alpha=\E w^2/(2\E\chi^3)$, smoothness and a Lipschitz bound for
$\alpha$, and bounds on mixed derivatives of $\Phi$ on compact subintervals.
The positive constant $m_\beta$ of \eqref{b:eq:main-analytic-constants} is
defined in \cref{cor:sc-mbeta}.

We use three external inputs: \cref{inp:ob-scalar}, on the regularity and
stability of the scalar equation; \cref{inp:full-rsb}, which states that the
support of $\mu_\beta$ is $[0,q_\beta]$; and the identity
\eqref{b:eq:identity-interval}, which states that $\Gamma(t)=t$ for
$0<t<q_\beta$. The derivative bounds are used in
\cref{sec:flow,sec:cubic,sec:fractional}. The stationarity identities enter
the coefficient computation of \cref{sec:cubic}, and the constant $m_\beta$
enters the cubic pressure deficit there and, through it, the choice of
parameters in \cref{sec:fractional,sec:references,sec:barrier-proof}.

The statements of this section are known, some of them in a slightly
different form. We include proofs because the later sections need them with
explicit uniformity. For each fixed distribution function, bounds on all
positive spatial derivatives, and their convergence when the profile is
approximated, were proved by Auffinger and Chen
\cite[Propositions~1--2]{auffinger-chen-2015} (see also
\cite[Proposition~2.4]{full-rsb}). The bounds of order at most three, the
martingale identity \eqref{b:eq:u-martingale}, and the expectation of
\eqref{b:eq:chi-evolution} are proved in \cite[Lemma~3.1]{oai-barriers}. Our
proof of \cref{lem:sc-derivatives} extends the Feynman--Kac argument given
there to every order. The identities \eqref{b:eq:os-jet-generator} with
$j\le3$ and the first two identities of \eqref{eq:sc-moments} appear in
\cite[Lemma~2.5, (2.12)--(2.16)]{full-rsb}; the formulas for $\Gamma'$
given there are stochastic versions of the differentiation formulas of
\cite[Proposition~3]{auffinger-chen-2015}. The identity $\Gamma(q)=q$ at
points of the support goes back to Talagrand
\cite{talagrand-2006-parisi-measures} (see
\cite[Theorem~5 and its proof]{auffinger-chen-2015}); it is part of the
characterization of Parisi measures by Chen \cite[Proposition~1]{chen-2017}
and by Jagannath and Tobasco \cite[Proposition~1.1]{jagannath-tobasco-2017},
and we use it in the form proved in \cite[Proposition~3.2]{oai-barriers}.
Since $\supp\mu_\beta=[0,q_\beta]$ \cite[Theorem~1.1]{full-rsb}, the identity holds on
all of $(0,q_\beta)$. In the proof
of \cite[Theorem~2(ii)]{auffinger-chen-2015}, Auffinger and Chen derive the
quotient formula for $\alpha$ and use an induction to prove its smoothness.
The same quotient argument, including the passage from almost every point
to every point, was carried out near $t=0$ in
\cite[proof of Lemma~7.2, Step~3, (7.8)--(7.12)]{quantitative-parisi}, in a
rescaled normalization. The smoothness of $\alpha_\beta$ on
$(0,q_\beta)$ is also part of \cref{inp:full-rsb} (see
\cref{cor:inp-identity}). The use of the second stationarity identity in the
covariance expansion of \cref{sec:cubic} is new.
% CHECK: Auffinger--Chen arXiv:1303.3573v1 read: Propositions 1--3 (Sections 2--3),
% Theorem 5 (Section 3; its proof credits Talagrand, Parisi measures, Prop. 3.2),
% Theorem 2(ii) with proof in Section 4.2, quotient (40); Lopatto arXiv:2607.11756v3, Proposition 2.4 and
% Lemma 2.5, (2.12)--(2.16), pp. 5--7; Aronow--Lopatto arXiv:2609.09103v1,
% proof of Lemma 7.2, Step 3, (7.8)--(7.12), pp. 63--64.

\paragraph{Notation.} Throughout this section, $\beta>1$ is fixed and
$b=\beta^2$. We use the conventions of \cref{b:sec:scalar}. A distribution
function is a nondecreasing, right-continuous function
$\alpha\colon[0,1]\to[0,1]$ with $\alpha(1)=1$; equivalently,
$\alpha(t)=\mu([0,t])$ for a probability measure $\mu$ on $[0,1]$. Let
$\alpha$ be a step distribution function, with breakpoints
$0=t_0<t_1<\dots<t_p=1$ and $\alpha=m_i$ on $[t_{i-1},t_i)$. Then
$\Phi_\alpha(1,z)=\log(2\cosh z)$, and the Gaussian recursion
\eqref{eq:inp-scalar-recursion} reads
\begin{equation}\label{eq:sc-recursion}
 \Phi_\alpha(t,z)=\frac1{m_i}\log\E\exp\Bigl(m_i\,
   \Phi_\alpha\bigl(t_i,z+\beta\sqrt{t_i-t}\,Z\bigr)\Bigr),
 \qquad t_{i-1}\le t<t_i,
\end{equation}
where $Z$ is a standard Gaussian variable, and where the right side is read
as $\E\Phi_\alpha(t_i,z+\beta\sqrt{t_i-t}\,Z)$ when $m_i=0$. For a
general distribution function $\alpha$, the function $\Phi_\alpha$ is the
uniform limit of $\Phi_{\alpha_N}$ along any sequence of step distribution
functions with $\norm{\alpha_N-\alpha}_{L^1}\to0$ (\cref{inp:ob-scalar}\,(i)). We
write
\[
 f_0=\Phi_\alpha,\qquad f_j=f_j^\alpha=\partial_z^j\Phi_\alpha\quad(j\ge1),
\]
so that $f_1=u$, $f_2=\chi$, and $f_3=w$. For $(t,z)\in[0,1]\times\R$, let
$(Y^{t,z}_r)_{t\le r\le1}$ be the solution of
\[
 \dd Y^{t,z}_r=\beta\dd B_r+b\,\alpha(r)\,u(r,Y^{t,z}_r)\dd r,\qquad
 Y^{t,z}_t=z,
\]
where $B$ is a standard Brownian motion. This equation has a unique strong
solution (\cref{inp:ob-scalar}\,(iii) and (v)), and the scalar diffusion of
\cref{sec:inputs} is $Y=Y^{0,0}$. We write $\norm{\cdot}_\infty$ for the
supremum over $[0,1]\times\R$, or over $\R$ for a function of $z$ alone.
From \cref{inp:ob-scalar}, we use the bounds $\abs u\le1$ and $0<\chi\le1$
(the proof of \cite[Lemma~3.1]{oai-barriers} gives the sharper bound
$\chi\le1-u^2$, which we do not need), the joint continuity of $\Phi_\alpha$,
$u$, and $\chi$, the well-posedness of the scalar diffusion, and the
stability estimate in part (ii) there,
\[
 \sup_{t,z}\bigl(\abs{\Phi_\alpha-\Phi_{\tilde\alpha}}
   +\abs{u_\alpha-u_{\tilde\alpha}}+\abs{\chi_\alpha-\chi_{\tilde\alpha}}\bigr)
 \le C_\beta\norm{\alpha-\tilde\alpha}_{L^1},
\]
which holds for all distribution functions $\alpha$ and $\tilde\alpha$.

\section{Bounds on spatial derivatives}\label{sec:sc-derivatives}

The derivative bounds below are uniform over all distribution functions. We
need this uniformity because the later sections apply the bounds to step
approximations and to profiles other than $\alpha_\beta$, such as
$\max\{\alpha_\beta,3s\}$ with $0\le s\le1/3$. The constants are defined by $C_1=C_2=1$ and, for
$j\ge3$,
\begin{equation}\label{eq:sc-constants}
 C_j=e^{jb}\Bigl(\norm{\partial_z^j\log(2\cosh z)}_\infty
   +\frac b2\sum_{a=2}^{j-2}\binom ja C_{a+1}C_{j+1-a}\Bigr).
\end{equation}
The sum is empty for $j=3$, and it involves only $C_3,\dots,C_{j-1}$. Every
derivative of positive order of $\log(2\cosh z)$ is a polynomial in
$\tanh z$ and is bounded.

\begin{lemma}[Bounds on spatial derivatives]\label{lem:sc-derivatives}
Let $\alpha$ be a distribution function. For all $j\ge1$, the derivative
$f_j=\partial_z^j\Phi_\alpha$ exists and is continuous on $[0,1]\times\R$,
and
\begin{equation}\label{eq:sc-derivative-bounds}
 \sup_{(t,z)\in[0,1]\times\R}\abs{\partial_z^j\Phi_\alpha(t,z)}\le C_j.
\end{equation}
Moreover, the following hold.
\begin{enumerate}[label=\textup{(\roman*)}]
\item If $\alpha_N$ are distribution functions such that
  $\norm{\alpha_N-\alpha}_{L^1}\to0$, then
  $\norm{\partial_z^j\Phi_{\alpha_N}-\partial_z^j\Phi_\alpha}_\infty\to0$ for
  all $j\ge1$.
\item For all $j\ge0$, $t\in[0,1]$, and $z\in\R$,
  \begin{equation}\label{eq:sc-integrated-pde}
   f_j(t,z)=\partial_z^j\log(2\cosh z)
     +\frac b2\int_t^1\bigl(f_{j+2}+\alpha\,\partial_z^j(u^2)\bigr)(r,z)\dd r,
  \end{equation}
  where
  $\partial_z^j(u^2)=\sum_{a=0}^j\binom ja f_{a+1}f_{j+1-a}$.
\end{enumerate}
\end{lemma}

The bounds concern derivatives of positive order only; the function
$\Phi_\alpha$ itself grows linearly in $z$. By \eqref{eq:sc-integrated-pde},
the Parisi equation and its spatial derivatives hold for almost every $t$,
and at every $t$ at which $\alpha$ is continuous.

\begin{proof}
We first recall why, for a fixed step distribution function, all spatial
derivatives exist and are bounded, with bounds that may depend on the steps.
We then write the equation for $f_j$ as a linear equation and use its
Feynman--Kac representation to obtain bounds independent of the steps.
Finally, we pass to general distribution functions by a forward difference
argument and prove \eqref{eq:sc-integrated-pde}.

\paragraph{Step 1: a fixed step distribution function.}
Let $\alpha$ be a step distribution function as in \eqref{eq:sc-recursion},
fix $i$, and let $\phi=\Phi_\alpha(t_i,\cdot)$. Suppose that $\phi$ is
smooth, that all its derivatives of positive order are bounded, and that
$\abs{\phi(z)}\le c(1+\abs z)$ for some $c$. These properties hold for
$\phi(z)=\log(2\cosh z)$, which is the case $i=p$. For
$t_{i-1}\le t\le t_i$, the function $\Phi_\alpha(t,\cdot)$ is $1/m_i$ times
the logarithm of the convolution of $e^{m_i\phi}$ with the centered Gaussian
density of variance $b(t_i-t)$, or the convolution of $\phi$ itself when
$m_i=0$. The linear growth of $\phi$ justifies differentiation under the
integral. We find that, for $j\ge1$, $f_j(t,z)$ is a polynomial, with
coefficients depending only on $j$ and $m_i$, in expectations of products of
$\phi'(z+\sigma Z),\dots,\phi^{(j)}(z+\sigma Z)$ under the probability
measure with density proportional to $e^{m_i\phi(z+\sigma Z)}$, where
$\sigma=\beta\sqrt{t_i-t}$. In particular, $f_j$ is bounded and
continuous on $[t_{i-1},t_i]\times\R$. For $t<t_i$, the Gaussian
convolution is smooth in $(t,z)$, and $e^{m_i\Phi_\alpha}$ (or
$\Phi_\alpha$ itself, when $m_i=0$) solves the backward heat equation
$-\partial_tv=(b/2)\partial_{zz}v$. It follows that
$\Phi_\alpha$ is smooth on $[t_{i-1},t_i)\times\R$, where it satisfies
\[
 -\partial_t\Phi_\alpha=\frac b2\bigl(\partial_{zz}\Phi_\alpha
   +m_i(\partial_z\Phi_\alpha)^2\bigr).
\]
The function $\Phi_\alpha(t_{i-1},\cdot)$ has the three properties assumed for
$\phi$ (it grows at most linearly because $\abs u\le1$). Proceeding backward
through the steps, we obtain these properties on every strip; this argument
is recorded in the proof of \cite[Lemma~3.1]{oai-barriers}. The bounds it gives depend on the number and
the location of the steps, and the next two steps remove this dependence.

\paragraph{Step 2: the differentiated equation.}
Differentiating the equation $j$ times in $z$ on each strip
$[t_{i-1},t_i)\times\R$, we find that
\begin{equation}\label{eq:sc-differentiated}
 -\partial_tf_j=\frac b2\bigl(f_{j+2}+\alpha\,\partial_z^j(u^2)\bigr)
 \qquad(j\ge0),
\end{equation}
with $\partial_z^j(u^2)$ as in \eqref{eq:sc-integrated-pde}. Let $j\ge3$. In
the sum defining $\partial_z^j(u^2)$, the terms with $a\in\{0,j\}$ add up to
$2uf_{j+1}$, and the terms with $a\in\{1,j-1\}$ add up to $2j\chi f_j$.
Then \eqref{eq:sc-differentiated} becomes
\begin{equation}\label{eq:sc-linear}
 -\partial_tf_j=\frac b2\partial_{zz}f_j+b\alpha u\,\partial_zf_j
   +jb\alpha\chi f_j+\mathcal Q_j,\qquad
 \mathcal Q_j=\frac{b\alpha}2\sum_{a=2}^{j-2}\binom ja f_{a+1}f_{j+1-a}.
\end{equation}
The sum defining $\mathcal Q_j$ is empty for $j=3$, and its factors are derivatives of
orders between $3$ and $j-1$. We regard \eqref{eq:sc-linear} as a linear
equation for $f_j$, with coefficients $u$ and $\chi$ and source term $\mathcal Q_j$.

\paragraph{Step 3: bounds independent of the steps.}
We prove \eqref{eq:sc-derivative-bounds} for step distribution functions by
induction on $j$. For $j\in\{1,2\}$, it is the bound $\abs u\le1$,
$0<\chi\le1$ of \cref{inp:ob-scalar}. Let $j\ge3$, and suppose that
$\norm{f_l}_\infty\le C_l$ for $1\le l\le j-1$. Since $0\le\alpha\le1$, we
have
\[
 \norm{\mathcal Q_j}_\infty\le\frac b2\sum_{a=2}^{j-2}\binom ja C_{a+1}C_{j+1-a}.
\]
Fix $(t,z)\in[0,1)\times\R$ and let $Y_r=Y^{t,z}_r$. We abbreviate
\[
 \Xi_r=jb\int_t^r\alpha(r')\chi(r',Y_{r'})\dd r',\qquad t\le r\le1,
\]
so that $0\le\Xi_r\le jb$. By \eqref{eq:sc-linear}, the drift of
$f_j(r,Y_r)$ in It\^o's formula is
\[
 \partial_rf_j+b\alpha u\,\partial_zf_j+\frac b2\partial_{zz}f_j
 =-jb\alpha\chi f_j-\mathcal Q_j,
\]
evaluated at $(r,Y_r)$. On each interval $[t_{i-1},t_i)\cap[t,1]$, we then
have
\[
 \dd\bigl(e^{\Xi_r}f_j(r,Y_r)\bigr)
 =-e^{\Xi_r}\mathcal Q_j(r,Y_r)\dd r
   +\beta e^{\Xi_r}f_{j+1}(r,Y_r)\dd B_r.
\]
Recall from Step~1 that the functions $f_j$ and $f_{j+1}$ are bounded and
continuous on $[0,1]\times\R$. The process
\[
 e^{\Xi_r}f_j(r,Y_r)
   +\int_t^re^{\Xi_{r'}}\mathcal Q_j(r',Y_{r'})\dd r',\qquad t\le r\le1,
\]
is therefore continuous, and it is a martingale. Comparing its expectations
at $r=t$ and $r=1$, we obtain
\[
 f_j(t,z)=\E\Bigl[e^{\Xi_1}\,\partial_z^j\log(2\cosh z)\big|_{z=Y_1}
   +\int_t^1e^{\Xi_r}\mathcal Q_j(r,Y_r)\dd r\Bigr].
\]
Since $0\le\Xi_r\le jb$ for $t\le r\le1$, we have
\[
 \abs{f_j(t,z)}\le e^{jb}\bigl(\norm{\partial_z^j\log(2\cosh z)}_\infty
   +\norm{\mathcal Q_j}_\infty\bigr)\le C_j.
\]
The same bound holds at $t=1$, which completes the induction. The
elapsed time is at most one, and the only property of $\alpha$ we used is
$0\le\alpha\le1$, so the constants depend neither on the number
nor on the location of the steps.

\paragraph{Step 4: general distribution functions.}
Let $\alpha$ be a distribution function, and let $\alpha_N$ be step
distribution functions such that $\norm{\alpha_N-\alpha}_{L^1}\to0$; for
instance, $\alpha_N(t)=\alpha((k+1)/N)$ for $k/N\le t<(k+1)/N$ and
$\alpha_N(1)=1$, which converges to $\alpha$ pointwise by right continuity.
We write $f_j^N=f_j^{\alpha_N}$. The stability estimate of
\cref{inp:ob-scalar} implies that $\Phi_{\alpha_N}$, $f_1^N$, and $f_2^N$
converge uniformly to $\Phi_\alpha$, $u$, and $\chi$. Let $j\ge3$, and suppose that
$f_{j-1}^N$ converges uniformly to $f_{j-1}=\partial_z^{j-1}\Phi_\alpha$,
which is continuous. For $\zeta>0$, Taylor's theorem and Step~3 imply
that
\[
 \Bigl|f_j^N(t,z)-\frac{f_{j-1}^N(t,z+\zeta)-f_{j-1}^N(t,z)}\zeta\Bigr|
 \le\frac{C_{j+1}\zeta}2,
\]
and that
\[
 \norm{f_j^N-f_j^M}_\infty
 \le\frac2\zeta\norm{f_{j-1}^N-f_{j-1}^M}_\infty+C_{j+1}\zeta.
\]
Given $\eps>0$, we first choose $\zeta$ with $C_{j+1}\zeta<\eps/2$, and then
$N_0$
such that the first term is less than $\eps/2$ for $N,M\ge N_0$. Then
$f_j^N$ converges uniformly to a continuous function $\widehat f_j$ with
$\norm{\widehat f_j}_\infty\le C_j$. Letting $N\to\infty$ in
$f_{j-1}^N(t,z)=f_{j-1}^N(t,0)+\int_0^zf_j^N(t,y)\dd y$, we find that
$\partial_zf_{j-1}$ exists and equals $\widehat f_j$. By induction on $j$, this proves
the existence and continuity of $f_j$, the bound
\eqref{eq:sc-derivative-bounds}, and the uniform convergence
$f_j^N\to f_j$ for all $j\ge1$. Now that \eqref{eq:sc-derivative-bounds}
holds for all distribution functions, the same argument applies to any
sequence of distribution functions $\alpha_N$ with
$\norm{\alpha_N-\alpha}_{L^1}\to0$, step functions or not. This proves (i).

\paragraph{Step 5: the integrated equation.}
Let $\alpha$ be a step distribution function and $j\ge0$. For fixed $z$, the
function $t\mapsto f_j(t,z)$ is continuous on $[0,1]$ by Step~1, and
continuously differentiable on each interval $[t_{i-1},t_i)$, with derivative
given by \eqref{eq:sc-differentiated}. Integrating over the steps, we obtain
\eqref{eq:sc-integrated-pde}. For a general $\alpha$, we choose $\alpha_N$ as
in Step~4 and write $u_N=f_1^N$. Then $\Phi_{\alpha_N}\to\Phi_\alpha$ and
$f_l^N\to f_l$ uniformly for all $l\ge1$, and
\[
 \Bigl|\int_t^1\bigl(\alpha_N\partial_z^j(u_N^2)-\alpha\,\partial_z^j(u^2)\bigr)
   (r,z)\dd r\Bigr|
 \le\norm{\alpha_N-\alpha}_{L^1}\norm{\partial_z^j(u_N^2)}_\infty
   +\norm{\partial_z^j(u_N^2)-\partial_z^j(u^2)}_\infty.
\]
The right side tends to zero by \eqref{eq:sc-derivative-bounds} and the
uniform convergence. Letting $N\to\infty$ in \eqref{eq:sc-integrated-pde}
for $\alpha_N$ yields \eqref{eq:sc-integrated-pde} for $\alpha$, which
completes the proof.
\end{proof}

\section{Stochastic identities along the scalar diffusion}
\label{sec:sc-ito}

We now record the equations satisfied by the spatial derivatives of
$\Phi_\alpha$ along the scalar diffusion. As usual, a differential identity
such as \eqref{b:eq:os-jet-generator} below means that the corresponding
integrated identity holds almost surely, simultaneously for all
$0\le s\le t\le1$. For a polynomial $\Theta$ in the variables
$x_1,\dots,x_k$, we define the polynomials in $x_1,\dots,x_{k+1}$
\begin{equation}\label{eq:sc-generator}
 \Theta^0=\frac b2\sum_{i,l=1}^kx_{i+1}x_{l+1}\,\partial_i\partial_l\Theta,
 \qquad
 \Theta^1=-b\sum_{i=1}^k\partial_i\Theta
   \sum_{a=1}^{i-1}\binom{i-1}ax_{a+1}x_{i+1-a},
\end{equation}
where $\partial_i$ is the derivative in $x_i$. We write
$\Theta_t=\Theta(f_1(t,Y_t),\dots,f_k(t,Y_t))$, and $\Theta^0_t$ and
$\Theta^1_t$ for the analogous expressions with $k+1$ arguments.

\begin{lemma}[Stochastic identities along the scalar diffusion]
\label{lem:sc-ito}
Let $\alpha$ be a distribution function, and let $Y=Y^{0,0}$. Write
$u_t=u(t,Y_t)$, $\chi_t=\chi(t,Y_t)$, and $w_t=w(t,Y_t)$.
\begin{enumerate}[label=\textup{(\roman*)}]
\item For all $j\ge1$,
  \begin{equation}\label{b:eq:os-jet-generator}
   \dd f_j(t,Y_t)=\beta f_{j+1}(t,Y_t)\dd B_t
     -b\,\alpha(t)\sum_{a=1}^{j-1}\binom{j-1}a
       f_{a+1}(t,Y_t)f_{j+1-a}(t,Y_t)\dd t.
  \end{equation}
  In particular,
  \begin{align}
   \dd u_t&=\beta\chi_t\dd B_t,\label{b:eq:u-martingale}\\
   \dd\chi_t&=\beta w_t\dd B_t-b\,\alpha(t)\chi_t^2\dd t.
     \label{b:eq:chi-evolution}
  \end{align}
\item Let $\Theta$ be a polynomial in $x_1,\dots,x_k$. Then
  $t\mapsto\E\Theta_t$ is Lipschitz on $[0,1]$, and for
  $0\le s\le t\le1$,
  \begin{equation}\label{eq:sc-moment-identity}
   \E\Theta_t-\E\Theta_s
   =\int_s^t\bigl(\E\Theta^0_r+\alpha(r)\E\Theta^1_r\bigr)\dd r.
  \end{equation}
\item For $0\le s\le t\le1$,
  \begin{equation}\label{eq:sc-moments}
  \begin{aligned}
   \Gamma(t)-\Gamma(s)&=b\int_s^t\E\chi_r^2\dd r,\\
   \E\chi_t^2-\E\chi_s^2
     &=b\int_s^t\bigl(\E w_r^2-2\alpha(r)\E\chi_r^3\bigr)\dd r,\\
   \E w_t^2-\E w_s^2
     &=b\int_s^t\E\bigl(f_4(r,Y_r)^2-6\alpha(r)\chi_rw_r^2\bigr)\dd r,\\
   \E\chi_t^3-\E\chi_s^3
     &=3b\int_s^t\E\bigl(\chi_rw_r^2-\alpha(r)\chi_r^4\bigr)\dd r.
  \end{aligned}
  \end{equation}
  In particular, $\Gamma$ is continuously differentiable with
  $\Gamma'(t)=b\E\chi_t^2$. The function $t\mapsto\E\chi_t^2$ has a right
  derivative at every $t\in[0,1)$, and a derivative at every point of
  continuity of $\alpha$, given by
  \begin{equation}\label{b:eq:contact-derivative}
   \frac{\dd}{\dd t}\E\chi_t^2=b\bigl(\E w_t^2-2\alpha(t)\E\chi_t^3\bigr).
  \end{equation}
\end{enumerate}
\end{lemma}

Parts (i) and (ii) hold, with the same proof, for the diffusion $Y^{t,z}$
started from any $(t,z)\in[0,1]\times\R$, on the time interval $[t,1]$.

\begin{proof}
We prove (i) first for step distribution functions, for which the equation
holds classically between the breakpoints, and then pass to the limit along
a coupling of the diffusions. Parts (ii) and (iii) then follow from It\^o's
formula.

\paragraph{Step 1: step distribution functions.}
Let $\alpha$ be a step distribution function. By Step~1 of the proof of
\cref{lem:sc-derivatives}, each $f_j$ is smooth on the strips
$[t_{i-1},t_i)\times\R$, bounded, and continuous on $[0,1]\times\R$. On each
strip, It\^o's formula shows that the martingale part of $f_j(t,Y_t)$ is
$\beta f_{j+1}(t,Y_t)\dd B_t$. Using \eqref{eq:sc-differentiated}, we find
that its drift is
\[
 \partial_tf_j+b\alpha u f_{j+1}+\frac b2f_{j+2}
 =-\frac{b\alpha}2\sum_{a=0}^j\binom ja f_{a+1}f_{j+1-a}+b\alpha uf_{j+1}
 =-\frac{b\alpha}2\sum_{a=1}^{j-1}\binom ja f_{a+1}f_{j+1-a},
\]
evaluated at $(t,Y_t)$. In the second equality, we used that the terms with
$a\in\{0,j\}$ equal $2uf_{j+1}$. By Pascal's rule
$\binom ja=\binom{j-1}a+\binom{j-1}{a-1}$ and the substitution
$a\mapsto j-a$ in the sum containing $\binom{j-1}{a-1}$, we have
\[
 \sum_{a=1}^{j-1}\binom ja f_{a+1}f_{j+1-a}
 =2\sum_{a=1}^{j-1}\binom{j-1}a f_{a+1}f_{j+1-a}.
\]
This proves \eqref{b:eq:os-jet-generator} on each interval $[t_{i-1},t_i)$.
The process $f_j(t,Y_t)$ is continuous in $t$, and the integrated identities
on the successive intervals add up to \eqref{b:eq:os-jet-generator} on
$[0,1]$.

\paragraph{Step 2: general distribution functions.}
Let $\alpha$ be a distribution function, let $\alpha_N$ be step distribution
functions such that $\norm{\alpha_N-\alpha}_{L^1}\to0$, and let $Y^N$ be the
scalar diffusion of $\alpha_N$, driven by the same Brownian motion $B$ as
$Y$. Let $u_N=f_1^{\alpha_N}$. Recall that $\abs{u_N}\le1$, and that $u$ is
$1$-Lipschitz in $z$ because $\abs\chi\le1$. Then
\[
 \abs{Y^N_t-Y_t}\le b\int_0^t\abs{\alpha_N(r)-\alpha(r)}\dd r
   +b\norm{u_N-u}_\infty+b\int_0^t\abs{Y^N_r-Y_r}\dd r ,
\]
and Gronwall's inequality yields
\[
 \sup_{0\le t\le1}\abs{Y^N_t-Y_t}
 \le be^b\bigl(\norm{\alpha_N-\alpha}_{L^1}+\norm{u_N-u}_\infty\bigr)
\]
for every realization of $B$. The right side tends to zero by the stability
estimate of \cref{inp:ob-scalar}. Using \cref{lem:sc-derivatives}, we find
that, for each $l\ge1$,
\[
 \abs{f_l^{\alpha_N}(t,Y^N_t)-f_l(t,Y_t)}
 \le\norm{f_l^{\alpha_N}-f_l}_\infty+C_{l+1}\abs{Y^N_t-Y_t},
\]
which tends to zero uniformly in $t$ and in the realization. We now let
$N\to\infty$ in \eqref{b:eq:os-jet-generator} for $\alpha_N$, in its
integrated form between fixed times $s\le t$. The terms $f_j(t,Y_t)$ and
$f_j(s,Y_s)$ converge uniformly. The drift integrals also converge
uniformly, because their integrands are bounded by
\cref{lem:sc-derivatives} and converge uniformly, and
$\norm{\alpha_N-\alpha}_{L^1}\to0$. By It\^o's isometry, the stochastic
integrals converge in $L^2$. We conclude that \eqref{b:eq:os-jet-generator}
holds almost surely for each fixed $s\le t$. Since both sides are continuous
in $(s,t)$, it holds almost surely for all $s\le t$ simultaneously. The cases
$j=1$ and $j=2$ are \eqref{b:eq:u-martingale} and
\eqref{b:eq:chi-evolution}.

\paragraph{Step 3: polynomial moments.}
It follows from (i) that the processes $f_1(t,Y_t),\dots,f_{k+1}(t,Y_t)$ are
continuous semimartingales with
$\dd\langle f_i(\cdot,Y),f_l(\cdot,Y)\rangle_t
 =b\,f_{i+1}(t,Y_t)f_{l+1}(t,Y_t)\dd t$. By It\^o's formula, we have
\[
 \dd\Theta_t=\bigl(\Theta^0_t+\alpha(t)\Theta^1_t\bigr)\dd t
   +\beta\sum_{i=1}^k\partial_i\Theta\bigl(f_1(t,Y_t),\dots,f_k(t,Y_t)\bigr)
     f_{i+1}(t,Y_t)\dd B_t.
\]
Here the second-order term of It\^o's formula produces $\Theta^0_t$, and the
drifts in \eqref{b:eq:os-jet-generator} produce $\alpha(t)\Theta^1_t$. By
\cref{lem:sc-derivatives}, the integrand of the stochastic integral is
bounded. Taking expectations, we obtain \eqref{eq:sc-moment-identity}. The
integrand in \eqref{eq:sc-moment-identity} is bounded by the supremum of
$\abs{\Theta^0}+\abs{\Theta^1}$ over the box $\prod_{i\le k+1}[-C_i,C_i]$,
which proves the Lipschitz bound.

\paragraph{Step 4: four moments.}
For $\Theta=x_1^2$, $x_2^2$, $x_3^2$, and $x_2^3$, the definition
\eqref{eq:sc-generator} gives
\[
\begin{aligned}
 (x_1^2)^0&=bx_2^2,&(x_1^2)^1&=0,\\
 (x_2^2)^0&=bx_3^2,&(x_2^2)^1&=-2bx_2^3,\\
 (x_3^2)^0&=bx_4^2,&(x_3^2)^1&=-6bx_2x_3^2,\\
 (x_2^3)^0&=3bx_2x_3^2,&(x_2^3)^1&=-3bx_2^4.
\end{aligned}
\]
For example, for $\Theta=x_3^2$ we have $\Theta^0=(b/2)\cdot2x_4^2$ and
$\Theta^1=-b\cdot2x_3\bigl(\binom21x_2x_3+\binom22x_3x_2\bigr)$. Inserting
these into \eqref{eq:sc-moment-identity}, we obtain \eqref{eq:sc-moments}.
The integrand $\E\chi_r^2$ in the first identity is continuous in $r$ by
(ii), and the first identity implies that $\Gamma$ is continuously
differentiable with $\Gamma'=b\E\chi^2$. In the second identity, the
functions $r\mapsto\E w_r^2$ and $r\mapsto\E\chi_r^3$ are continuous by (ii),
and $\alpha$ is right-continuous. The integrand is therefore right-continuous,
and it is continuous at every point of continuity of $\alpha$. The
fundamental theorem of calculus now gives \eqref{b:eq:contact-derivative}.
\end{proof}

\section{The Parisi distribution function on \texorpdfstring{$(0,q_\beta)$}{(0,q)}}
\label{sec:sc-support}

From now on, $\alpha=\alpha_\beta$ is the Parisi distribution function, which
is unique by \cref{inp:unique}, $\mu_\beta$ is the Parisi measure, $Y$ is the
associated scalar diffusion, and $u_t$, $\chi_t$, $w_t$ are as in
\cref{lem:sc-ito}. By \cref{inp:full-rsb}, $\supp\mu_\beta=[0,q_\beta]$ with
$0<q_\beta<1$, and $\mu_\beta$ has an atom at $q_\beta$. We set
\begin{equation}\label{eq:sc-lipschitz-constant}
 K_\beta=\frac{\beta^5}2\bigl(C_4^2+12C_3^2+6\bigr),
\end{equation}
with $C_3$ and $C_4$ as in \eqref{eq:sc-constants}.

\begin{lemma}[The Parisi distribution function on the support]
\label{b:lem:os-scalar-jets}
The following hold.
\begin{enumerate}[label=\textup{(\roman*)}]
\item For all $t\in(0,q_\beta)$,
  % NOTE: second label via the kernel \label saved by amsmath (amsmath forbids
  % two \label's in one display); cleveref also records b:eq:os-contacts@cref.
  \begin{equation}\label{b:eq:contacts}\csname ltx@label\endcsname{b:eq:os-contacts}
   \beta^2\E\chi(t,Y_t)^2=1,\qquad
   \E w(t,Y_t)^2=2\alpha(t)\E\chi(t,Y_t)^3.
  \end{equation}
\item For all $t\in(0,q_\beta)$, we have $\E\chi_t^3\ge\beta^{-3}$ and
  \begin{equation}\label{eq:sc-quotient}
   \alpha(t)=\frac{\E w_t^2}{2\E\chi_t^3}.
  \end{equation}
\item We have $\alpha(t)>0$ for all $t\in(0,1]$. For
  $0<s\le t<q_\beta$,
  \[
   0\le\alpha(t)-\alpha(s)=\mu_\beta((s,t])\le K_\beta(t-s).
  \]
  In particular, $\mu_\beta$ has no atoms in $(0,q_\beta)$.
\item The function $\alpha$ is infinitely differentiable on $(0,q_\beta)$.
  For every polynomial $\Theta$, the function $t\mapsto\E\Theta_t$ of
  \cref{lem:sc-ito} is infinitely differentiable on $(0,q_\beta)$. In
  particular, $t\mapsto\E\chi_t^3$ and
  $t\mapsto\alpha(t)^2\beta^6\E\chi_t^3$ are infinitely differentiable on
  $(0,q_\beta)$.
\item The function $\Phi_\alpha$ is infinitely differentiable on
  $(0,q_\beta)\times\R$. For every compact interval $I\subset(0,q_\beta)$
  and all integers $i,j\ge0$ with $i+j\ge1$, there exists
  $C(\beta,I,i,j)<\infty$ such that
  \[
   \sup_{t\in I}\abs{\alpha^{(i)}(t)}
   +\sup_{(t,z)\in I\times\R}\abs{\partial_t^i\partial_z^j\Phi_\alpha(t,z)}
   \le C(\beta,I,i,j).
  \]
\end{enumerate}
\end{lemma}

The proof of (i) uses the identity $\Gamma(t)=t$, which is a property of
$\alpha_\beta$ itself. The step distribution functions used in the proof of
\cref{lem:sc-ito} do not satisfy it, and this is why we established the
identities of \cref{lem:sc-ito} for every distribution function first.

\begin{proof}
We first obtain the two identities in (i): the first at every point, from
$\Gamma(t)=t$, and the second at almost every point, from the constancy of
$\E\chi_t^2$. Since $\alpha$ is right-continuous, the second identity then
holds at every point, which gives the quotient formula
\eqref{eq:sc-quotient}. Using the quotient formula and the Lipschitz bounds of
\cref{lem:sc-ito}, we show that $\alpha$ is Lipschitz, and an induction on the
order of differentiability gives its smoothness. The
bounds on mixed derivatives then follow from the differentiated Parisi
equation.

\paragraph{Step 1: the first identity.}
The first identity is also stated in \cref{cor:inp-identity}; we recall the
argument. By \eqref{b:eq:identity-interval}, $\Gamma(t)=t$ for
$0<t<q_\beta$. By \cref{lem:sc-ito}\,(iii), $\Gamma$ is continuously
differentiable with $\Gamma'(t)=b\E\chi_t^2$. Differentiating, we find that
$b\E\chi_t^2=\Gamma'(t)=1$ for all $t\in(0,q_\beta)$, which is the first
identity in \eqref{b:eq:contacts}.

\paragraph{Step 2: the second identity almost everywhere.}
By Step~1, $\E\chi_t^2=\beta^{-2}$ is constant on $(0,q_\beta)$. The
second identity in \eqref{eq:sc-moments} implies that, for $0<s<t<q_\beta$,
\[
 0=\E\chi_t^2-\E\chi_s^2
 =b\int_s^t\bigl(\E w_r^2-2\alpha(r)\E\chi_r^3\bigr)\dd r.
\]
The integral of the bounded function
$r\mapsto\E w_r^2-2\alpha(r)\E\chi_r^3$ vanishes over every subinterval of
$(0,q_\beta)$. Lebesgue's differentiation theorem then implies that
$\E w_r^2=2\alpha(r)\E\chi_r^3$ for almost every $r\in(0,q_\beta)$.

\paragraph{Step 3: the second identity at every point, and the quotient.}
Let $t\in(0,q_\beta)$. By Step~2, there exist $r_n\downarrow t$ such that
$\E w_{r_n}^2=2\alpha(r_n)\E\chi_{r_n}^3$ for all $n$. Recall from \cref{lem:sc-ito}\,(ii) that the
functions $r\mapsto\E w_r^2$ and $r\mapsto\E\chi_r^3$ are continuous, and
that $\alpha$ is right-continuous. Letting $n\to\infty$, we find that
$\E w_t^2=2\alpha(t)\E\chi_t^3$. This proves (i). Since $\chi_t>0$ and
$x\mapsto x^{3/2}$ is convex on $[0,\infty)$, Jensen's inequality and Step~1
imply that
\[
 \E\chi_t^3=\E\bigl[(\chi_t^2)^{3/2}\bigr]\ge\bigl(\E\chi_t^2\bigr)^{3/2}
 =\beta^{-3}.
\]
Dividing the second identity in \eqref{b:eq:contacts} by $2\E\chi_t^3$, we
obtain \eqref{eq:sc-quotient}, which proves (ii).

\paragraph{Step 4: positivity and the Lipschitz bound.}
Since $0\in\supp\mu_\beta$, every interval $[0,t]$ with $t>0$ has
positive $\mu_\beta$-mass, so $\alpha(t)>0$ for $t\in(0,1]$. By
\cref{lem:sc-derivatives}, we have $0<\chi\le1$, $\abs w\le C_3$, and
$\abs{f_4}\le C_4$. Inserting these bounds into the last two identities in
\eqref{eq:sc-moments}, we find that, for all $s,t\in[0,1]$,
\begin{equation}\label{eq:sc-moment-lipschitz}
 \abs{\E w_t^2-\E w_s^2}\le b\bigl(C_4^2+6C_3^2\bigr)\abs{t-s},\qquad
 \abs{\E\chi_t^3-\E\chi_s^3}\le3b\bigl(C_3^2+1\bigr)\abs{t-s}.
\end{equation}
Let $0<s\le t<q_\beta$. By \eqref{eq:sc-quotient} at $t$ and the identity
$\E w_s^2=2\alpha(s)\E\chi_s^3$ at $s$, we have
\[
 \alpha(t)-\alpha(s)
 =\frac{\E w_t^2-\E w_s^2}{2\E\chi_t^3}
   -\alpha(s)\,\frac{\E\chi_t^3-\E\chi_s^3}{\E\chi_t^3}.
\]
Using \eqref{eq:sc-moment-lipschitz}, $\E\chi_t^3\ge\beta^{-3}$, and
$0\le\alpha(s)\le1$, we conclude that
\[
 \abs{\alpha(t)-\alpha(s)}
 \le\beta^3\Bigl(\frac{b}{2}\bigl(C_4^2+6C_3^2\bigr)
   +3b\bigl(C_3^2+1\bigr)\Bigr)(t-s)=K_\beta(t-s).
\]
Since $\alpha$ is nondecreasing and $\alpha(t)-\alpha(s)=\mu_\beta((s,t])$,
this proves (iii); the continuity of $\alpha$ on $(0,q_\beta)$ excludes
atoms there.

\paragraph{Step 5: smoothness of $\alpha$ and of the moments.}
Let $I_0=(0,q_\beta)$. We prove by induction on $i\ge0$ that
$\alpha\in C^i(I_0)$ and that $t\mapsto\E\Theta_t$ belongs to $C^{i+1}(I_0)$
for every polynomial $\Theta$. By Step~4, $\alpha$ is continuous on $I_0$.
Then the integrand $r\mapsto\E\Theta^0_r+\alpha(r)\E\Theta^1_r$ in
\eqref{eq:sc-moment-identity} is continuous on $I_0$, and $\E\Theta_t$ is
continuously differentiable on $I_0$ by the fundamental theorem of calculus,
with
\begin{equation}\label{eq:sc-moment-derivative}
 \frac{\dd}{\dd t}\E\Theta_t=\E\Theta^0_t+\alpha(t)\E\Theta^1_t
 \qquad(t\in I_0).
\end{equation}
This proves the claim for $i=0$. Suppose that it holds for some $i\ge0$.
Since the moments $\E w_t^2$ and $\E\chi_t^3$ belong to $C^{i+1}(I_0)$ and
$\E\chi_t^3\ge\beta^{-3}$, the formula \eqref{eq:sc-quotient} implies that
$\alpha\in C^{i+1}(I_0)$. By the induction hypothesis, the moments
$\E\Theta^0_t$ and $\E\Theta^1_t$ belong to $C^{i+1}(I_0)$. Then the right
side of \eqref{eq:sc-moment-derivative} belongs to $C^{i+1}(I_0)$, and we
obtain $\E\Theta_t\in C^{i+2}(I_0)$. This completes the induction and proves
(iv).

\paragraph{Step 6: mixed derivatives.}
Let $j\ge0$. By \eqref{eq:sc-integrated-pde} and the continuity of $\alpha$
on $I_0$, the function $t\mapsto f_j(t,z)$ is continuously differentiable on
$I_0$ for every $z$, and
\begin{equation}\label{eq:sc-time-derivative}
 \partial_tf_j=-\frac b2\bigl(f_{j+2}+\alpha\,\partial_z^j(u^2)\bigr)
 \qquad\text{on }I_0\times\R.
\end{equation}
The right side is a polynomial in $\alpha$ and $f_1,\dots,f_{j+2}$, and it is
continuous on $I_0\times\R$. We claim that, for every $i\ge1$, the
derivative $\partial_t^if_j$ exists on $I_0\times\R$ and is a polynomial,
with coefficients depending only on $b$, $i$, and $j$, in
$\alpha,\alpha',\dots,\alpha^{(i-1)}$ (evaluated at $t$) and
$f_1,\dots,f_{j+2i}$ (evaluated at $(t,z)$). For $i=1$, this is
\eqref{eq:sc-time-derivative}. If it holds for $i$, we differentiate the
polynomial in $t$, using (iv) for the derivatives of $\alpha$ and
\eqref{eq:sc-time-derivative} for the derivatives of $f_1,\dots,f_{j+2i}$,
which gives the claim for $i+1$. Let $I\subset I_0$ be compact. On
$I\times\R$, the functions $\alpha,\alpha',\dots,\alpha^{(i)}$ are bounded
because they are continuous on $I$, and we have $\abs{f_l}\le C_l$ by
\cref{lem:sc-derivatives}. We obtain the bound in (v), with a constant
depending on $\beta$, $I$, $i$, and $j$. Finally, the right side of \eqref{eq:sc-time-derivative} for $j+1$
is the $z$-derivative of the right side for $j$, because
$\partial_z\partial_z^j(u^2)=\partial_z^{j+1}(u^2)$. This shows that
$\partial_z\partial_tf_j=\partial_tf_{j+1}$, and by Schwarz's theorem and
induction on $i$, we have $\partial_z\partial_t^if_j=\partial_t^if_{j+1}$.
Every partial derivative of $\Phi_\alpha$, taken in any order, is therefore
one of the continuous functions $\partial_t^i\partial_z^j\Phi_\alpha$, and
$\Phi_\alpha$ is infinitely differentiable on $I_0\times\R$. This completes
the proof.
\end{proof}

\begin{remark}\label{rem:sc-upper-atom}
The second identity in \eqref{b:eq:contacts} does not extend to
$t=q_\beta$. Indeed, $t\mapsto\E w_t^2$ and $t\mapsto\E\chi_t^3$ are
continuous on $[0,1]$, and letting $t\uparrow q_\beta$ in
\eqref{b:eq:contacts}, we find that
$\E w_{q_\beta}^2=2\alpha(q_\beta-)\E\chi_{q_\beta}^3$. However,
$\E\chi_{q_\beta}^3>0$, and
$\alpha(q_\beta-)=1-\mu_\beta(\{q_\beta\})<1=\alpha(q_\beta)$ because
$\mu_\beta$ has an atom at $q_\beta$. The barrier argument uses
\eqref{b:eq:contacts} and \cref{b:lem:os-scalar-jets} only on compact
subintervals of $(0,q_\beta)$.
\end{remark}

We conclude with the positive constant that enters the cubic pressure
deficit. The quantity bounded below in \eqref{b:eq:main-analytic-constants} is
the scalar factor of the cubic term in the expansion of
\cref{b:prop:os-expansion}.

\begin{corollary}[The constant $m_\beta$]\label{cor:sc-mbeta}
Let $I_\beta\subset(0,q_\beta)$ be a nonempty compact interval, and set
$m_\beta=\beta^3\alpha(\min I_\beta)^2$. Then $m_\beta>0$, and
\begin{equation}\label{b:eq:main-analytic-constants}
 m_\beta\le\inf_{r\in I_\beta}\alpha(r)^2\beta^6\E\chi(r,Y_r)^3.
\end{equation}
\end{corollary}

\begin{proof}
By \cref{b:lem:os-scalar-jets}\,(iii), we have $\alpha(\min I_\beta)>0$, so
$m_\beta>0$. Let $r\in I_\beta$. Since $\alpha$ is nondecreasing, we have
$\alpha(r)\ge\alpha(\min I_\beta)$, and \cref{b:lem:os-scalar-jets}\,(ii)
gives $\E\chi(r,Y_r)^3\ge\beta^{-3}$. Multiplying these bounds, we find that
$\alpha(r)^2\beta^6\E\chi(r,Y_r)^3\ge\beta^3\alpha(\min I_\beta)^2=m_\beta$.
\end{proof}

The constant $m_\beta$ depends only on $\beta$ and $I_\beta$. In the proof of
\cref{thm:barrier-730}, the interval $I_\beta$ is fixed in \cref{sec:references},
before any of the small parameters of the local comparison is chosen.
